\centerline{\bf \Huge Муниципальный листочек}
\centerline{\bf \Huge Теория чисел + перекладывания треугольников}

\problem{1}
Вася получил за год несколько оценок по математике, всего их было меньше 100. Ровно треть из них - тройки, ровно четверть четвёрки, ровно пятая часть - пятёрки. А сколько Вася получил двоек? В задаче предполагается, что возможные оценки - это $5,4,3$ и 2 .

\problem{2}
У Пети в двух карманах было по одинаковому количеству монет. Он высыпал все эти монеты на стол и подсчитал, что орлов выпало на 7 больше, чем решек. Докажите, что он ошибся.


\problem{3}
Найдите три подряд идущих трехзначных числа таких, что каждое делится на квадрат своей последней цифры.


\textit{Напомним принцип перекладывания треугольников}

\textit{В треугольниках $A B C$ и $A^{\prime} B^{\prime} C^{\prime} A B=A^{\prime} B^{\prime}$ и $\angle B A C+\angle B^{\prime} A^{\prime} C^{\prime}= 180^{\circ}$. Тогда равенство сторон $B C$ и $B^{\prime} C^{\prime}$ равносильно равенству углов $\angle B C A$ и $\angle B^{\prime} C^{\prime} A^{\prime}$.}

\problem{4}
В выпуклом четырёхугольнике $A B C D$, в котором $A B=C D$, на сторонах $A B$ и $C D$ выбраны точки $K$ и $M$ соответственно. Оказалось, что $A M=K C, B M=K D$. Докажите, что угол между прямыми $A B$ и $K M$ равен углу между прямыми $K M$ и $C D$.

\problem{5}
В неравнобедренном треугольнике $A B C$ биссектрисы $A A_1$ и $B B_1$ пересекаются в точке $I$. Найдите $\angle C$, если $A_1 I=B_1 I$.

\problem{6}
В выпуклом четырёхугольнике $A B C D \angle A=\angle C$. На продолжении $B A$ за точкой $A$ выбрали точку $E$ такую, что $B C=A E$. Оказалось, что $\angle B D C=\angle A D E$ и $\angle A B D= \angle C A D$. Докажите, что $A C \perp B D$.